\makeatletter
\def\input@path{{./}{../}{../../}{../../../}{../../../../}{preamble/}{../preamble/}{../../preamble/}{../../../preamble/}}
\makeatother

\input{preamble_exams.tex}

\title{Grade 11 -- Term 1 \\ Classroom Activity: Experimental Probability}
\author{}
\date{}

\pagestyle{empty}

% ============================================================
% SOLUTION STAGE COMMAND
% ============================================================

\newcommand{\SolutionStage}[2]{%
	\Needspace{6\baselineskip}%
	\subsection{#2}%
}

\begin{document}

	% ============================================================
	% PROBLEM PAGES
	% ============================================================

	\newgeometry{
		left=0.55in,
		right=0.55in,
		top=0.40in,
		bottom=0.45in
	}

	\vspace*{-15mm}

	\begin{center}
		{\large\bfseries G11 T1 C6 -- Experimental Interval Probabilities}
	\end{center}

	\vspace{-3mm}

	Suppose a sample contains $n$ observations and exactly $k$ of them
	satisfy the stated condition. The experimental probability, or relative
	frequency, is
	\[
		\boxed{\widehat P(a\leq X\leq b)=\frac{k}{n}},
	\]
	where $n$ is the sample size and $k$ is the number of observations for
	which $a\leq X\leq b$. It may be reported in three equivalent forms:
	\[
		\widehat P=\frac{k}{n}
		=\text{decimal form}
		=\bigl(100\times\text{decimal form}\bigr)\%.
	\]
	The same counting method applies to conditions using $<$, $\leq$, $>$,
	or $\geq$. Pay close attention to whether each endpoint is included.

	\begin{enumerate}
		\item \textbf{Water-bottle filling lines (multiple samples).}
		Three filling lines produced the following sample fill volumes, in mL:
		\[
		\begin{array}{c|cccccccc}
			\text{Line A}&496&500&503&499&505&501&497&502\\
			\text{Line B}&494&498&500&506&501&503&499&502\\
			\text{Line C}&497&501&504&500&498&502&503&499
		\end{array}
		\]
		For the combined sample from all three lines, estimate
		$\widehat P(499\leq X\leq502)$. Give the answer as a fraction, a
		decimal, and a percentage.

		\item \textbf{Bike commute times (raw data).}
		The commute times, in minutes, for 15 rides were
		\[
			18,\ 22,\ 25,\ 19,\ 31,\ 24,\ 27,\ 20,\ 23,\ 29,\ 21,\ 26,\ 17,\ 24,\ 28.
		\]
		Estimate $\widehat P(X<24)$ as a fraction, a decimal, and a percentage.

		\item \textbf{Books borrowed (value--frequency table).}
		A librarian recorded the number of books borrowed by 40 visitors.
		\[
		\begin{array}{c|rrrrrrr}
			\text{Books }x&0&1&2&3&4&5&6\\ \hline
			\text{Frequency }f&3&7&10&9&6&4&1
		\end{array}
		\]
		Estimate $\widehat P(X\geq3)$ as a fraction, a decimal, and a percentage.

		\item \textbf{Parcel masses (class-interval table).}
		The masses of 50 parcels are grouped below. Each class includes its
		left endpoint but not its right endpoint.
		\[
		\begin{array}{c|ccccc}
			\text{Mass }x\text{ (kg)}&[0,2)&[2,4)&[4,6)&[6,8)&[8,10)\\ \hline
			\text{Frequency}&4&11&17&12&6
		\end{array}
		\]
		Estimate $\widehat P(2\leq X<8)$ as a fraction, a decimal, and a percentage.

		\item \textbf{Seedling growth (raw data).}
		The increases in height, in cm, of 20 seedlings were
		\[
		\begin{gathered}
		4.2,\ 5.0,\ 5.7,\ 6.1,\ 4.8,\ 5.5,\ 6.0,\ 6.4,\ 5.2,\ 4.5,\\
		5.9,\ 6.3,\ 5.0,\ 5.4,\ 4.9,\ 6.0,\ 5.8,\ 6.6,\ 5.1,\ 5.6.
		\end{gathered}
		\]
		Estimate $\widehat P(5.0<X\leq6.0)$ as a fraction, a decimal, and a percentage.

		\item \textbf{Customer service calls (value--frequency table).}
		A service center summarized call lengths, rounded to the nearest minute.
		\[
		\begin{array}{c|rrrrrrrr}
			\text{Length }x\text{ (min)}&2&3&4&5&6&7&8&9\\ \hline
			\text{Frequency }f&2&5&8&11&9&7&5&3
		\end{array}
		\]
		Estimate both (a) $\widehat P(X>6)$ and
		(b) $\widehat P(4\leq X<8)$. Give each answer in all three forms.

		\item \textbf{Daily solar-energy output (class-interval table).}
		The output recorded on 80 days is grouped below. Each class includes
		its left endpoint but not its right endpoint.
		\[
		\begin{array}{c|rrrrrr}
			\text{Output }x\text{ (kWh)}&[0,5)&[5,10)&[10,15)&[15,20)&[20,25)&[25,30)\\ \hline
			\text{Frequency}&5&12&21&24&13&5
		\end{array}
		\]
		Estimate both (a) $\widehat P(X\geq15)$ and
		(b) $\widehat P(10\leq X<25)$. Give each answer in all three forms.
	\end{enumerate}

% ============================================================
% SOLUTIONS
% ============================================================

\newpage
\restoregeometry

\vspace*{-15mm}

\section{Solutions}

\setlength{\columnseprule}{0.4pt}

\begin{multicols}{2}

	\subsection*{C6.1 -- Water-bottle filling lines}
	There are $8$ observations from each of three lines, so
	\[
		n=3(8)=24.
	\]
	The qualifying values are $500,499,501,502$ from line A;
	$500,501,499,502$ from line B; and $501,500,502,499$ from line C. Thus,
	\[
		k=4+4+4=12.
	\]
	Therefore,
	\[
		\widehat P(499\leq X\leq502)
		=\frac{k}{n}=\frac{12}{24}=\frac12=0.50=50\%.
	\]

	\subsection*{C6.2 -- Bike commute times}
	There are $n=15$ rides. The times below $24$ are
	$18,22,19,20,23,21,17$, so $k=7$. Therefore,
	\[
		\widehat P(X<24)
		=\frac{k}{n}=\frac{7}{15}\approx0.467\approx46.7\%.
	\]

	\subsection*{C6.3 -- Books borrowed}
	The sample size is the sum of all frequencies:
	\[
		n=3+7+10+9+6+4+1=40.
	\]
	For $X\geq3$, sum the frequencies at $3,4,5,$ and $6$:
	\[
		k=9+6+4+1=20.
	\]
	Therefore,
	\[
		\widehat P(X\geq3)
		=\frac{k}{n}=\frac{20}{40}=\frac12=0.50=50\%.
	\]

	\subsection*{C6.4 -- Parcel masses}
	The sample size is
	\[
		n=4+11+17+12+6=50.
	\]
	The included buckets are $[2,4)$, $[4,6)$, and $[6,8)$. Hence,
	\[
		k=11+17+12=40.
	\]
	Therefore,
	\[
		\widehat P(2\leq X<8)
		=\frac{k}{n}=\frac{40}{50}=\frac45=0.80=80\%.
	\]

	\subsection*{C6.5 -- Seedling growth}
	There are $n=20$ seedlings. The values satisfying $5.0<X\leq6.0$ are
	\[
		5.7,5.5,6.0,5.2,5.9,5.4,6.0,5.8,5.1,5.6,
	\]
	so $k=10$. Notice that the two values equal to $5.0$ are excluded, while
	the two values equal to $6.0$ are included. Therefore,
	\[
		\widehat P(5.0<X\leq6.0)
		=\frac{k}{n}=\frac{10}{20}=\frac12=0.50=50\%.
	\]

	\subsection*{C6.6 -- Customer service calls}
	For both parts, the sample size is
	\[
		n=2+5+8+11+9+7+5+3=50.
	\]
	\textbf{(a)} For $X>6$, sum the frequencies at $7,8,$ and $9$:
	\[
		k=7+5+3=15.
	\]
	Thus,
	\[
		\widehat P(X>6)=\frac{k}{n}=\frac{15}{50}=\frac3{10}=0.30=30\%.
	\]
	\textbf{(b)} For $4\leq X<8$, sum the frequencies at $4,5,6,$ and $7$:
	\[
		k=8+11+9+7=35.
	\]
	Thus,
	\[
		\widehat P(4\leq X<8)=\frac{k}{n}=\frac{35}{50}=\frac7{10}=0.70=70\%.
	\]

	\subsection*{C6.7 -- Daily solar-energy output}
	For both parts, the sample size is
	\[
		n=5+12+21+24+13+5=80.
	\]
	\textbf{(a)} The included buckets are $[15,20)$, $[20,25)$, and
	$[25,30)$. Therefore,
	\[
		k=24+13+5=42,
	\]
	and
	\[
		\widehat P(X\geq15)=\frac{k}{n}=\frac{42}{80}=\frac{21}{40}
		=0.525=52.5\%.
	\]
	\textbf{(b)} The included buckets are $[10,15)$, $[15,20)$, and
	$[20,25)$. Therefore,
	\[
		k=21+24+13=58,
	\]
	and
	\[
		\widehat P(10\leq X<25)=\frac{k}{n}=\frac{58}{80}=\frac{29}{40}
		=0.725=72.5\%.
	\]

\end{multicols}

\end{document}
